\section{The strongest objections}\label{sec:objections}
\begin{enumerate}
\item \textbf{CBT has a large evidence base.} Yes (Section~\ref{sec:evidence}).
The argument concerns mechanism and moderators, not average efficacy.
\item \textbf{CBT for adult ADHD works.} It beats matched controls in medicated adults with a skills-heavy protocol; the evidence is graded low to moderate (Section~\ref{sec:adhdcbt}).
\item \textbf{Patients practise when calm.} Accepted as a plausible defence (Section~\ref{sec:ingredients}).
Rehearsal could improve cognitive skill as well as behavioural routines; it does not identify habit formation as the mechanism.
\item \textbf{Exposure is behavioural.} This is why CBT cannot be equated with cognitive restructuring.
The distinction between a procedure and its mechanism matters (Section~\ref{sec:ingredients}).
\item \textbf{Mindfulness is not better.} The comparative evidence does not establish general superiority.
The hypothesis concerns when particular components help, not whether one therapy family should replace another.
\item \textbf{Cognitive therapy sometimes wins.} Yes, in social phobia \citep{ougrin2011exposure} and in one PTSD trial \citep{bryant2008cr}.
\item \textbf{The neuroscience is about acute stress.} Correct, and the human reappraisal-under-stress findings are mixed (Section~\ref{sec:stress}).
\item \textbf{Executive function does not predict CBT outcome.} Largely true for overall outcome (Section~\ref{sec:moderators}); the argument survives only in the narrower, skill-level form.
\item \textbf{Therapists scaffold the executive load.} Plausible and untested; a therapist who generates the evidence list lowers demand and could erase a moderator effect.
\item \textbf{Evolutionary age does not rank therapies.} Correct.
Conservation motivates a learning-based hypothesis but cannot establish comparative clinical benefit (Section~\ref{sec:learning}).
\item \textbf{Cognition can help before experience.} Correct.
The learning and timing evidence is reviewed in Section~\ref{sec:learning}; acceptance-based goals are distinguished from eliminating distress in Section~\ref{sec:clinicalformulation}.
\item \textbf{Bias is not proof that efficacy is illusory.} Correct.
Researcher allegiance, publication selection and reviewer preference are distinct; they apply to competing orientations too and do not establish treatment equivalence (Section~\ref{sec:selection}).
\end{enumerate}
\textbf{Limitations.} This is a narrative review by a non-clinician, with no independent duplicate screening.
Several key studies involve older adults, children or elevated-symptom rather than diagnosed-disorder samples.
Evidence from depression informs some anxiety questions, and laboratory learning is not a direct test of psychotherapy.
The clinical heuristic is not a validated taxonomy, and the added literature searches were targeted rather than exhaustive.
No mechanism proposed here should be treated as an established explanation of an individual patient's outcome.
